\abstractPT  % Do NOT modify this line

A U!REKA é uma aliança europeia de instituições de ensino superior cuja missão é preparar cidadãos para liderar a transição para cidades neutras em carbono e sustentáveis. Esta transição exige competências digitais e cívicas --- definidas pelo quadro DigComp~2.2 --- que a maioria das plataformas de aprendizagem digital não desenvolve com eficácia, em parte pela taxa de abandono que as caracteriza.

Para endereçar este problema, foi desenvolvido e avaliado o \emph{Play4Change}: uma plataforma móvel gamificada assente em três pilares. O primeiro é um motor de geração de conteúdo pedagógico alimentado por um modelo de linguagem (Mistral~AI), capaz de ingerir documentos e URLs e de activar percursos adaptativos personalizados quando o utilizador erra. O segundo é um sistema de distintivos com limiar de mestria que assinala cada competência adquirida. O terceiro engloba mecanismos de retenção --- sequências diárias de utilização (\emph{streaks}) e um modo de explicação holística por IA --- que combatem o abandono.

A avaliação preliminar com 12~utilizadores obteve uma pontuação SUS de 76/100 (banda ``Bom''), uma taxa de conclusão de tópicos de 59\% e uma resolução de 98,2\% das sessões adaptativas. Os dados revelaram uma clivagem bimodal: utilizadores com elevada familiaridade digital (SUS médio~89,4) adoptaram a plataforma com entusiasmo, enquanto utilizadores menos familiarizados (SUS médio~47,5) abandonaram cedo --- evidenciando que a barreira não é o conteúdo, mas o meio. O ciclo de \emph{feedback} por IA foi espontaneamente identificado como a funcionalidade mais valorizada por 7~dos 11~respondentes.

\vspace{0.5cm}
\noindent\textbf{Palavras-chave:} aprendizagem adaptativa, educação cívica, gamificação, inteligência artificial, DigComp, cidades sustentáveis
